% GENERATED FILE - do not hand-edit.
% Source of truth: data/calculus-reviewed/t-6.json in the zenith-website repo.
% Regenerate with: python3 scripts/build_calculus_tex.py
\documentclass[11pt]{article}
\usepackage[utf8]{inputenc}
\usepackage[T1]{fontenc}
\usepackage{amsmath,amssymb}
\usepackage{graphicx}
\usepackage[margin=1in]{geometry}
\graphicspath{{../..//}}
\setlength{\parindent}{0pt}
\setlength{\parskip}{0.6em}
\begin{document}
\section*{Chapter 6: Definite Integrals --- Test T6}

\bigskip\noindent\textbf{T1.}\ $\displaystyle\int_{-1}^{2} (x^2-2x-3)\,dx =$
\begin{list}{}{\setlength{\leftmargin}{3em}\setlength{\labelwidth}{2.5em}\setlength{\itemsep}{0pt}\setlength{\topsep}{2pt}}
\item[(A)] $-3$
\item[(B)] $0$
\item[(C)] $-9$
\item[(D)] $9$
\end{list}

\bigskip\noindent\textbf{T2.}\ $\displaystyle\int_{1}^{2} \dfrac{4x-1}{4x}\,dx =$
\begin{list}{}{\setlength{\leftmargin}{3em}\setlength{\labelwidth}{2.5em}\setlength{\itemsep}{0pt}\setlength{\topsep}{2pt}}
\item[(A)] $\dfrac{3}{4}$
\item[(B)] $1-\dfrac{1}{4}\ln 2$
\item[(C)] $1-\ln 2$
\item[(D)] $-\dfrac{1}{4}\ln 2$
\end{list}

\bigskip\noindent\textbf{T3.}\ $\displaystyle\int_{0}^{8} \dfrac{dt}{\sqrt{9-t}} =$
\begin{list}{}{\setlength{\leftmargin}{3em}\setlength{\labelwidth}{2.5em}\setlength{\itemsep}{0pt}\setlength{\topsep}{2pt}}
\item[(A)] $-2$
\item[(B)] $4$
\item[(C)] $-4$
\item[(D)] $8$
\end{list}

\bigskip\noindent\textbf{T4.}\ $\displaystyle\int_{-1}^{0} \sqrt{8u+9}\,du =$
\begin{list}{}{\setlength{\leftmargin}{3em}\setlength{\labelwidth}{2.5em}\setlength{\itemsep}{0pt}\setlength{\topsep}{2pt}}
\item[(A)] $2$
\item[(B)] $\dfrac{13}{6}$
\item[(C)] $\dfrac{26}{3}$
\item[(D)] $\dfrac{13}{3}$
\end{list}

\bigskip\noindent\textbf{T5.}\ $\displaystyle\int_{2}^{3} \dfrac{dy}{3y-4} =$
\begin{list}{}{\setlength{\leftmargin}{3em}\setlength{\labelwidth}{2.5em}\setlength{\itemsep}{0pt}\setlength{\topsep}{2pt}}
\item[(A)] $\ln 5$
\item[(B)] $\dfrac{1}{3}\ln\dfrac{5}{2}$
\item[(C)] $\dfrac{25}{16}$
\item[(D)] $\ln\sqrt[3]{5}$
\end{list}

\bigskip\noindent\textbf{T6.}\ $\displaystyle\int_{0}^{\sqrt{12}} \dfrac{x\,dx}{\sqrt{16-x^2}} =$
\begin{list}{}{\setlength{\leftmargin}{3em}\setlength{\labelwidth}{2.5em}\setlength{\itemsep}{0pt}\setlength{\topsep}{2pt}}
\item[(A)] $1$
\item[(B)] $\dfrac{\pi}{6}$
\item[(C)] $-1$
\item[(D)] $2$
\end{list}

\bigskip\noindent\textbf{T7.}\ $\displaystyle\int_{0}^{1} (3t-1)^3\,dt =$
\begin{list}{}{\setlength{\leftmargin}{3em}\setlength{\labelwidth}{2.5em}\setlength{\itemsep}{0pt}\setlength{\topsep}{2pt}}
\item[(A)] $\dfrac{1}{4}$
\item[(B)] $\dfrac{5}{4}$
\item[(C)] $0$
\item[(D)] $4$
\end{list}

\bigskip\noindent\textbf{T8.}\ $\displaystyle\int_{4}^{9} \dfrac{3+x}{2\sqrt{x}}\,dx =$
\begin{list}{}{\setlength{\leftmargin}{3em}\setlength{\labelwidth}{2.5em}\setlength{\itemsep}{0pt}\setlength{\topsep}{2pt}}
\item[(A)] $\dfrac{28}{3}$
\item[(B)] $\dfrac{41}{3}$
\item[(C)] $\dfrac{100}{3}$
\item[(D)] $\dfrac{5}{3}$
\end{list}

\bigskip\noindent\textbf{T9.}\ $\displaystyle\int_{-4}^{4} \dfrac{dx}{16+x^2} =$
\begin{list}{}{\setlength{\leftmargin}{3em}\setlength{\labelwidth}{2.5em}\setlength{\itemsep}{0pt}\setlength{\topsep}{2pt}}
\item[(A)] $\dfrac{\pi}{4}$
\item[(B)] $0$
\item[(C)] $\dfrac{\pi}{8}$
\item[(D)] $-\dfrac{\pi}{4}$
\end{list}

\bigskip\noindent\textbf{T10.}\ $\displaystyle\int_{0}^{2} e^{-x}\,dx =$
\begin{list}{}{\setlength{\leftmargin}{3em}\setlength{\labelwidth}{2.5em}\setlength{\itemsep}{0pt}\setlength{\topsep}{2pt}}
\item[(A)] $\dfrac{1}{e^2}-1$
\item[(B)] $-\dfrac{1}{e^2}$
\item[(C)] $1-\dfrac{1}{e^2}$
\item[(D)] $\dfrac{1}{e^2}$
\end{list}

\bigskip\noindent\textbf{T11.}\ $\displaystyle\int_{0}^{2} xe^{x^2}\,dx =$
\begin{list}{}{\setlength{\leftmargin}{3em}\setlength{\labelwidth}{2.5em}\setlength{\itemsep}{0pt}\setlength{\topsep}{2pt}}
\item[(A)] $e^4-1$
\item[(B)] $\dfrac{1}{2}(e^4-1)$
\item[(C)] $2(e^4-1)$
\item[(D)] $\dfrac{e^4}{2}-1$
\end{list}

\bigskip\noindent\textbf{T12.}\ $\displaystyle\int_{0}^{\pi/6} \sin 3\theta\,d\theta =$
\begin{list}{}{\setlength{\leftmargin}{3em}\setlength{\labelwidth}{2.5em}\setlength{\itemsep}{0pt}\setlength{\topsep}{2pt}}
\item[(A)] $3$
\item[(B)] $\dfrac{1}{3}$
\item[(C)] $-1$
\item[(D)] $-\dfrac{1}{3}$
\end{list}

\bigskip\noindent\textbf{T13.}\ $\displaystyle\int_{1}^{2} \dfrac{dz}{4-z} =$
\begin{list}{}{\setlength{\leftmargin}{3em}\setlength{\labelwidth}{2.5em}\setlength{\itemsep}{0pt}\setlength{\topsep}{2pt}}
\item[(A)] $-\ln\dfrac{3}{2}$
\item[(B)] $\ln 3$
\item[(C)] $\dfrac{1}{2}\ln\dfrac{3}{2}$
\item[(D)] $\ln\dfrac{3}{2}$
\end{list}

\bigskip\noindent\textbf{T14.}\ If we let $x=4\sin\theta$, then $\displaystyle\int_{2}^{4} \dfrac{\sqrt{16-x^2}}{x}\,dx$ is equivalent to
\begin{list}{}{\setlength{\leftmargin}{3em}\setlength{\labelwidth}{2.5em}\setlength{\itemsep}{0pt}\setlength{\topsep}{2pt}}
\item[(A)] $\displaystyle 2\int_{\pi/6}^{\pi/2} \dfrac{\cos^2\theta}{\sin\theta}\,d\theta$
\item[(B)] $\displaystyle \int_{\pi/6}^{\pi/2} \dfrac{\cos\theta}{\sin\theta}\,d\theta$
\item[(C)] $\displaystyle 4\int_{\pi/6}^{\pi/2} \dfrac{\cos^2\theta}{\sin\theta}\,d\theta$
\item[(D)] $\displaystyle \int_{2}^{4} \dfrac{\cos\theta}{\sin\theta}\,d\theta$
\end{list}

\bigskip\noindent\textbf{T15.}\ $\displaystyle\int_{0}^{\pi/2} \cos^2\theta\sin\theta\,d\theta =$
\begin{list}{}{\setlength{\leftmargin}{3em}\setlength{\labelwidth}{2.5em}\setlength{\itemsep}{0pt}\setlength{\topsep}{2pt}}
\item[(A)] $-\dfrac{2}{3}$
\item[(B)] $1$
\item[(C)] $\dfrac{2}{3}$
\item[(D)] $\dfrac{1}{3}$
\end{list}

\bigskip\noindent\textbf{T16.}\ $\displaystyle\int_{1}^{e^2} \dfrac{\ln x}{x}\,dx =$
\begin{list}{}{\setlength{\leftmargin}{3em}\setlength{\labelwidth}{2.5em}\setlength{\itemsep}{0pt}\setlength{\topsep}{2pt}}
\item[(A)] $2$
\item[(B)] $\dfrac{1}{2}(e^2-1)$
\item[(C)] $4$
\item[(D)] $e^2-1$
\end{list}

\bigskip\noindent\textbf{T17.}\ $\displaystyle\int_{0}^{2} xe^x\,dx =$
\begin{list}{}{\setlength{\leftmargin}{3em}\setlength{\labelwidth}{2.5em}\setlength{\itemsep}{0pt}\setlength{\topsep}{2pt}}
\item[(A)] $-1$
\item[(B)] $e^2+1$
\item[(C)] $1$
\item[(D)] $e^2-1$
\end{list}

\bigskip\noindent\textbf{T18.}\ $\displaystyle\int_{0}^{\pi/6} \dfrac{\cos\theta}{1+3\sin\theta}\,d\theta =$
\begin{list}{}{\setlength{\leftmargin}{3em}\setlength{\labelwidth}{2.5em}\setlength{\itemsep}{0pt}\setlength{\topsep}{2pt}}
\item[(A)] $\ln\dfrac{5}{2}$
\item[(B)] $\dfrac{3}{8}$
\item[(C)] $-\dfrac{1}{3}\ln\dfrac{5}{2}$
\item[(D)] $\ln\sqrt[3]{5/2}$
\end{list}

\bigskip\noindent\textbf{T19.}\ $\displaystyle\int_{\sqrt{2}}^{3} \dfrac{u\,du}{u^2-1} =$
\begin{list}{}{\setlength{\leftmargin}{3em}\setlength{\labelwidth}{2.5em}\setlength{\itemsep}{0pt}\setlength{\topsep}{2pt}}
\item[(A)] $\ln\sqrt{8}$
\item[(B)] $\dfrac{8}{9}$
\item[(C)] $\ln\dfrac{3}{2}$
\item[(D)] $\ln 3$
\end{list}

\bigskip\noindent\textbf{T20.}\ $\displaystyle\int_{\sqrt{2}}^{3} \dfrac{u\,du}{(u^2-1)^2} =$
\begin{list}{}{\setlength{\leftmargin}{3em}\setlength{\labelwidth}{2.5em}\setlength{\itemsep}{0pt}\setlength{\topsep}{2pt}}
\item[(A)] $-\dfrac{1}{3}$
\item[(B)] $-\dfrac{7}{16}$
\item[(C)] $\dfrac{7}{16}$
\item[(D)] $\dfrac{1}{3}$
\end{list}

\bigskip\noindent\textbf{T21.}\ $\displaystyle\int_{\pi/18}^{\pi/6} \dfrac{\cos 3x\,dx}{\sin^2 3x} =$
\begin{list}{}{\setlength{\leftmargin}{3em}\setlength{\labelwidth}{2.5em}\setlength{\itemsep}{0pt}\setlength{\topsep}{2pt}}
\item[(A)] $1$
\item[(B)] $\dfrac{1}{3}$
\item[(C)] $-\dfrac{1}{3}$
\item[(D)] $-1$
\end{list}

\bigskip\noindent\textbf{T22.}\ $\displaystyle\int_{0}^{1} \dfrac{e^{-x}+2}{e^{-x}}\,dx =$
\begin{list}{}{\setlength{\leftmargin}{3em}\setlength{\labelwidth}{2.5em}\setlength{\itemsep}{0pt}\setlength{\topsep}{2pt}}
\item[(A)] $2e-1$
\item[(B)] $2+e$
\item[(C)] $\dfrac{1}{e}$
\item[(D)] $1+e$
\end{list}

\bigskip\noindent\textbf{T23.}\ $\displaystyle\int_{0}^{1} \dfrac{e^x}{e^x+2}\,dx =$
\begin{list}{}{\setlength{\leftmargin}{3em}\setlength{\labelwidth}{2.5em}\setlength{\itemsep}{0pt}\setlength{\topsep}{2pt}}
\item[(A)] $\ln 2$
\item[(B)] $1+e$
\item[(C)] $-\ln 2$
\item[(D)] $\ln\dfrac{e+2}{3}$
\end{list}

\bigskip\noindent\textbf{T24.}\ If we let $x=\tan\theta$, then $\displaystyle\int_{0}^{1} \sqrt{1+x^2}\,dx$ is equivalent to
\begin{list}{}{\setlength{\leftmargin}{3em}\setlength{\labelwidth}{2.5em}\setlength{\itemsep}{0pt}\setlength{\topsep}{2pt}}
\item[(A)] $\displaystyle\int_{0}^{\pi/4} \sec\theta\,d\theta$
\item[(B)] $\displaystyle\int_{0}^{1} \sec^3\theta\,d\theta$
\item[(C)] $\displaystyle\int_{0}^{\pi/4} \sec^3\theta\,d\theta$
\item[(D)] $\displaystyle\int_{0}^{\pi/4} \sec^2\theta\tan\theta\,d\theta$
\end{list}

\bigskip\noindent\textbf{T25.}\ If the substitution $u=\sqrt{x+1}$ is used, then $\displaystyle\int_{0}^{8} \dfrac{dx}{x\sqrt{x+1}}$ is equivalent to
\begin{list}{}{\setlength{\leftmargin}{3em}\setlength{\labelwidth}{2.5em}\setlength{\itemsep}{0pt}\setlength{\topsep}{2pt}}
\item[(A)] $\displaystyle\int_{1}^{3} \dfrac{du}{u^2-1}$
\item[(B)] $\displaystyle\int_{1}^{3} \dfrac{2\,du}{u^2-1}$
\item[(C)] $\displaystyle 2\int_{0}^{3} \dfrac{du}{(u-1)(u+1)}$
\item[(D)] $\displaystyle 2\int_{1}^{3} \dfrac{du}{u(u^2-1)}$
\end{list}

\bigskip\noindent\textbf{T26.}\ The table lists values of $f$ and its derivative: $\begin{array}{c|rrrrr}x&0&2&4&6&8\\\hline f(x)&9&14&18&13&6\\f'(x)&4&3&-2&-4&0\end{array}$. Evaluate $\int_8^0f'(x)\,dx$.
\begin{list}{}{\setlength{\leftmargin}{3em}\setlength{\labelwidth}{2.5em}\setlength{\itemsep}{0pt}\setlength{\topsep}{2pt}}
\item[(A)] $-\dfrac{1}{2}$
\item[(B)] $3$
\item[(C)] $4$
\item[(D)] $-1$
\end{list}

\bigskip\noindent\textbf{T27.}\ Using a midpoint Riemann Sum with 4 equal width subintervals, find the approximate area of the region under $y=8x-x^2$ on $[0,8]$.
\begin{list}{}{\setlength{\leftmargin}{3em}\setlength{\labelwidth}{2.5em}\setlength{\itemsep}{0pt}\setlength{\topsep}{2pt}}
\item[(A)] $44$
\item[(B)] $72$
\item[(C)] $88$
\item[(D)] $128$
\end{list}

\bigskip\noindent\textbf{T28.}\ The graph of a continuous function $f$ passes through the points $(3,2)$, $(5,8)$, $(6,6)$, and $(9,12)$. Using trapezoids, we estimate that $\displaystyle\int_{3}^{9} f(x)\,dx \approx$
\begin{list}{}{\setlength{\leftmargin}{3em}\setlength{\labelwidth}{2.5em}\setlength{\itemsep}{0pt}\setlength{\topsep}{2pt}}
\item[(A)] $40$
\item[(B)] $44$
\item[(C)] $48$
\item[(D)] $52$
\end{list}

\bigskip\noindent\textbf{T29.}\ Using left and right sums with 3 equal subintervals on $[1,4]$ for $y=1/x$, which gives the resulting lower and upper bounds directly for $\ln4$?
\begin{list}{}{\setlength{\leftmargin}{3em}\setlength{\labelwidth}{2.5em}\setlength{\itemsep}{0pt}\setlength{\topsep}{2pt}}
\item[(A)] $\dfrac{1}{4} < \displaystyle\int_{1}^{4} \dfrac{1}{x}\,dx < 1$
\item[(B)] $\dfrac{13}{12} < \displaystyle\int_{1}^{4} \dfrac{1}{x}\,dx < \dfrac{11}{6}$
\item[(C)] $\dfrac{1}{4} < \displaystyle\int_{3}^{4} \dfrac{1}{x}\,dx < \dfrac{1}{3}$
\item[(D)] $\dfrac{11}{6} < \displaystyle\int_{1}^{4} \dfrac{1}{x}\,dx < \dfrac{13}{12}$
\end{list}

\bigskip\noindent\textbf{T30.}\ Let $A = \displaystyle\int_{0}^{1} \sqrt{1-x^2}\,dx$. We estimate $A$ using the $L$, $R$, and $T$ approximations with $n=100$ subintervals. Which is true?
\begin{list}{}{\setlength{\leftmargin}{3em}\setlength{\labelwidth}{2.5em}\setlength{\itemsep}{0pt}\setlength{\topsep}{2pt}}
\item[(A)] $R < T < A < L$
\item[(B)] $L < T < A < R$
\item[(C)] $R < A < T < L$
\item[(D)] $L < A < T < R$
\end{list}

\bigskip\noindent\textbf{T31.}\ $\displaystyle\int_{-2}^{3} |x|\,dx =$
\begin{list}{}{\setlength{\leftmargin}{3em}\setlength{\labelwidth}{2.5em}\setlength{\itemsep}{0pt}\setlength{\topsep}{2pt}}
\item[(A)] $\dfrac{7}{2}$
\item[(B)] $4$
\item[(C)] $\dfrac{9}{2}$
\item[(D)] $\dfrac{13}{2}$
\end{list}

\bigskip\noindent\textbf{T32.}\ $\displaystyle\int_{-4}^{3} |x+2|\,dx =$
\begin{list}{}{\setlength{\leftmargin}{3em}\setlength{\labelwidth}{2.5em}\setlength{\itemsep}{0pt}\setlength{\topsep}{2pt}}
\item[(A)] $\dfrac{21}{2}$
\item[(B)] $\dfrac{25}{2}$
\item[(C)] $\dfrac{27}{2}$
\item[(D)] $\dfrac{29}{2}$
\end{list}

\bigskip\noindent\textbf{T33.}\ The average value of $y=\sqrt{36-x^2}$ on its domain is
\begin{list}{}{\setlength{\leftmargin}{3em}\setlength{\labelwidth}{2.5em}\setlength{\itemsep}{0pt}\setlength{\topsep}{2pt}}
\item[(A)] $3$
\item[(B)] $6$
\item[(C)] $\dfrac{3\pi}{2}$
\item[(D)] $3\pi$
\end{list}

\bigskip\noindent\textbf{T34.}\ The average value of $\sin x$ over the interval $\dfrac{\pi}{6} \leq x \leq \dfrac{\pi}{3}$ is
\begin{list}{}{\setlength{\leftmargin}{3em}\setlength{\labelwidth}{2.5em}\setlength{\itemsep}{0pt}\setlength{\topsep}{2pt}}
\item[(A)] $\dfrac{3}{\pi}$
\item[(B)] $\dfrac{3(\sqrt{3}-1)}{\pi}$
\item[(C)] $\dfrac{3}{2\pi}$
\item[(D)] $\dfrac{2}{3\pi}$
\end{list}

\bigskip\noindent\textbf{T35.}\ The average value of $\sec^2 x$ over the interval from $x=\dfrac{\pi}{6}$ to $x=\dfrac{\pi}{4}$ is
\begin{list}{}{\setlength{\leftmargin}{3em}\setlength{\labelwidth}{2.5em}\setlength{\itemsep}{0pt}\setlength{\topsep}{2pt}}
\item[(A)] $\dfrac{3\sqrt{3}}{\pi}$
\item[(B)] $\dfrac{\sqrt{3}}{\pi}$
\item[(C)] $\dfrac{4(3-\sqrt{3})}{\pi}$
\item[(D)] $\dfrac{3(3-\sqrt{3})}{\pi}$
\end{list}

\bigskip\noindent\textbf{T36.}\ Choose the Riemann Sum whose limit is the integral $\displaystyle\int_{2}^{6} x^3\,dx$.
\begin{list}{}{\setlength{\leftmargin}{3em}\setlength{\labelwidth}{2.5em}\setlength{\itemsep}{0pt}\setlength{\topsep}{2pt}}
\item[(A)] $\displaystyle\lim_{n\to\infty}\sum_{k=1}^n \left(\left(2+\dfrac{k}{n}\right)^3\cdot\left(\dfrac{4}{n}\right)\right)$
\item[(B)] $\displaystyle\lim_{n\to\infty}\sum_{k=1}^n \left(\left(2+\dfrac{4k}{n}\right)^3\cdot\left(\dfrac{1}{n}\right)\right)$
\item[(C)] $\displaystyle\lim_{n\to\infty}\sum_{k=1}^n \left(\left(2+\dfrac{4k}{n}\right)^3\cdot\left(\dfrac{4}{n}\right)\right)$
\item[(D)] $\displaystyle\lim_{n\to\infty}\sum_{k=1}^n \left(\left(2+\dfrac{k}{n}\right)^3\cdot\left(\dfrac{1}{n}\right)\right)$
\end{list}

\bigskip\noindent\textbf{T37.}\ Which expression is obtained by constructing the right-endpoint Riemann sum for $\int_0^{2\pi}\sin(4x)\,dx$ using $n$ equal subintervals of $[0,2\pi]$, then taking its limit?
\begin{list}{}{\setlength{\leftmargin}{3em}\setlength{\labelwidth}{2.5em}\setlength{\itemsep}{0pt}\setlength{\topsep}{2pt}}
\item[(A)] $\displaystyle\lim_{n\to\infty}\sum_{k=1}^n \left(\sin\left(\dfrac{4k}{n}\right)\cdot\left(\dfrac{1}{n}\right)\right)$
\item[(B)] $\displaystyle\lim_{n\to\infty}\sum_{k=1}^n \left(\sin\left(\dfrac{8\pi k}{n}\right)\cdot\left(\dfrac{1}{n}\right)\right)$
\item[(C)] $\displaystyle\lim_{n\to\infty}\sum_{k=1}^n \left(\sin\left(\dfrac{4k}{n}\right)\cdot\left(\dfrac{2\pi}{n}\right)\right)$
\item[(D)] $\displaystyle\lim_{n\to\infty}\sum_{k=1}^n \left(\sin\left(\dfrac{8\pi k}{n}\right)\cdot\left(\dfrac{2\pi}{n}\right)\right)$
\end{list}

\bigskip\noindent\textbf{T38.}\ Choose the integral that is the limit of the Riemann Sum
\[
\lim_{n\to\infty}\sum_{k=1}^n \left(\sin\left(2+\dfrac{8k}{n}\right)\cdot\left(\dfrac{4}{n}\right)\right).
\]
\begin{list}{}{\setlength{\leftmargin}{3em}\setlength{\labelwidth}{2.5em}\setlength{\itemsep}{0pt}\setlength{\topsep}{2pt}}
\item[(A)] $\displaystyle\int_{2}^{6} \sin(8x)\,dx$
\item[(B)] $\displaystyle\int_{0}^{4} \sin(2x+2)\,dx$
\item[(C)] $\displaystyle\int_{2}^{6} \sin(2x)\,dx$
\item[(D)] $\displaystyle\int_{0}^{4} \sin(8x+2)\,dx$
\end{list}

\bigskip\noindent\textbf{T39.}\ $\displaystyle\int_{\pi/6}^{\pi/3} \sin^3\alpha\cos\alpha\,d\alpha$ is equal to
\begin{list}{}{\setlength{\leftmargin}{3em}\setlength{\labelwidth}{2.5em}\setlength{\itemsep}{0pt}\setlength{\topsep}{2pt}}
\item[(A)] $\dfrac{3}{16}$
\item[(B)] $\dfrac{1}{8}$
\item[(C)] $-\dfrac{1}{8}$
\item[(D)] $-\dfrac{3}{16}$
\end{list}

\bigskip\noindent\textbf{TB1.}\ Find the average value on $[0,6]$ of the piecewise linear function joining $(0,0),(3,12),(6,12)$.
\begin{list}{}{\setlength{\leftmargin}{3em}\setlength{\labelwidth}{2.5em}\setlength{\itemsep}{0pt}\setlength{\topsep}{2pt}}
\item[(A)] $6$
\item[(B)] $8$
\item[(C)] $9$
\item[(D)] $12$
\end{list}

\bigskip\noindent\textbf{TB2.}\ The integral $\displaystyle\int_{-5}^{5} \sqrt{25-x^2}\,dx$ gives the area of
\begin{list}{}{\setlength{\leftmargin}{3em}\setlength{\labelwidth}{2.5em}\setlength{\itemsep}{0pt}\setlength{\topsep}{2pt}}
\item[(A)] a circle of radius 5
\item[(B)] a semicircle of radius 5
\item[(C)] a quadrant of a circle of radius 5
\item[(D)] half of an ellipse
\end{list}

\bigskip\noindent\textbf{TB3.}\ The graph of $f$ consists of line segments joining $(0,2),(3,8),(9,0)$. Which intervals contain a point $c$ where $f(c)$ equals its average on $[0,9]$? I. $[0,3]$; II. $[3,6]$; III. $[6,9]$.
\begin{list}{}{\setlength{\leftmargin}{3em}\setlength{\labelwidth}{2.5em}\setlength{\itemsep}{0pt}\setlength{\topsep}{2pt}}
\item[(A)] I only
\item[(B)] II only
\item[(C)] III only
\item[(D)] I and II only
\end{list}

\bigskip\noindent\textbf{TB4.}\ The graph of $f$ consists of line segments joining $(0,2),(3,8),(9,0)$. $\displaystyle\int_{0}^{2} f'(3x)\,dx =$
\begin{list}{}{\setlength{\leftmargin}{3em}\setlength{\labelwidth}{2.5em}\setlength{\itemsep}{0pt}\setlength{\topsep}{2pt}}
\item[(A)] $-2$
\item[(B)] $-\dfrac{2}{3}$
\item[(C)] $\dfrac{2}{3}$
\item[(D)] $2$
\end{list}

\bigskip\noindent\textbf{TB5.}\ The graph of $f$ consists of line segments joining $(0,2),(3,8),(9,0)$. Let $g(x) = \displaystyle\int_{0}^{2x} f(t)\,dt$; then $g'(1) =$
\begin{list}{}{\setlength{\leftmargin}{3em}\setlength{\labelwidth}{2.5em}\setlength{\itemsep}{0pt}\setlength{\topsep}{2pt}}
\item[(A)] $6$
\item[(B)] $8$
\item[(C)] $10$
\item[(D)] $12$
\end{list}

\bigskip\noindent\textbf{TB6.}\ The graph of $f$ consists of line segments joining $(0,2),(3,8),(9,0)$. Let $h(x) = x^2 - f(x)$. Find $\displaystyle\int_{0}^{9} h(x)\,dx$.
\begin{list}{}{\setlength{\leftmargin}{3em}\setlength{\labelwidth}{2.5em}\setlength{\itemsep}{0pt}\setlength{\topsep}{2pt}}
\item[(A)] $165$
\item[(B)] $204$
\item[(C)] $235$
\item[(D)] $243$
\end{list}

\bigskip\noindent\textbf{TB7.}\ If $f(x)$ is continuous on the closed interval $[a,b]$, then there exists at least one number $c$, $a<c<b$, such that $\displaystyle\int_{a}^{b} f(x)\,dx$ is equal to
\begin{list}{}{\setlength{\leftmargin}{3em}\setlength{\labelwidth}{2.5em}\setlength{\itemsep}{0pt}\setlength{\topsep}{2pt}}
\item[(A)] $\dfrac{f(c)}{b-a}$
\item[(B)] $f'(c)(b-a)$
\item[(C)] $f(c)(b-a)$
\item[(D)] $\dfrac{f'(c)}{b-a}$
\end{list}

\bigskip\noindent\textbf{TB8.}\ If $f(x)$ is continuous on the closed interval $[a,b]$ and $k$ is a constant, then $\displaystyle\int_{a}^{b} kf(x)\,dx$ is equal to
\begin{list}{}{\setlength{\leftmargin}{3em}\setlength{\labelwidth}{2.5em}\setlength{\itemsep}{0pt}\setlength{\topsep}{2pt}}
\item[(A)] $k(b-a)$
\item[(B)] $k[f(b)-f(a)]$
\item[(C)] $\dfrac{[kf(x)]^2}{2}\Big|_a^b$
\item[(D)] $k\displaystyle\int_{a}^{b} f(x)\,dx$
\end{list}

\bigskip\noindent\textbf{TB9.}\ $\displaystyle\dfrac{d}{dt}\int_{0}^{t} \sqrt{x^3+1}\,dx =$
\begin{list}{}{\setlength{\leftmargin}{3em}\setlength{\labelwidth}{2.5em}\setlength{\itemsep}{0pt}\setlength{\topsep}{2pt}}
\item[(A)] $\sqrt{t^3+1}$
\item[(B)] $\dfrac{\sqrt{t^3+1}}{3t^2}$
\item[(C)] $3x^2\sqrt{x^3+1}$
\item[(D)] $\sqrt{t^3+1}-1$
\end{list}

\bigskip\noindent\textbf{TB10.}\ If $F(u) = \displaystyle\int_{1}^{u} (3-x^2)^4\,dx$, then $F'(u)$ is equal to
\begin{list}{}{\setlength{\leftmargin}{3em}\setlength{\labelwidth}{2.5em}\setlength{\itemsep}{0pt}\setlength{\topsep}{2pt}}
\item[(A)] $-8u(3-u^2)^3$
\item[(B)] $(3-u^2)^4-1$
\item[(C)] $(3-u^2)^4$
\item[(D)] $-2u(3-u^2)^4$
\end{list}

\bigskip\noindent\textbf{TB11.}\ $\displaystyle\dfrac{d}{dx}\int_{\pi/2}^{x^3} \sqrt{\sin t}\,dt =$
\begin{list}{}{\setlength{\leftmargin}{3em}\setlength{\labelwidth}{2.5em}\setlength{\itemsep}{0pt}\setlength{\topsep}{2pt}}
\item[(A)] $\sqrt{\sin t^3}$
\item[(B)] $3x^2\sqrt{\sin x^3}-1$
\item[(C)] $\sqrt{\sin x^3}-1$
\item[(D)] $3x^2\sqrt{\sin x^3}$
\end{list}

\bigskip\noindent\textbf{TB12.}\ On the upper ellipse branch $x=5\cos\theta$, $y=4\sin\theta$, $0\leq\theta\leq\pi/3$, the integral $\int_{2.5}^5xy\,dx$ is equivalent to
\begin{list}{}{\setlength{\leftmargin}{3em}\setlength{\labelwidth}{2.5em}\setlength{\itemsep}{0pt}\setlength{\topsep}{2pt}}
\item[(A)] $\displaystyle 100\int_{\pi/3}^{0} \sin\theta\cos^2\theta\,d\theta$
\item[(B)] $\displaystyle 100\int_{2.5}^{5} \sin^2\theta\cos\theta\,d\theta$
\item[(C)] $\displaystyle 100\int_{0}^{\pi/3} \sin\theta\cos^2\theta\,d\theta$
\item[(D)] $\displaystyle 100\int_{0}^{\pi/3} \sin^2\theta\cos\theta\,d\theta$
\end{list}

\bigskip\noindent\textbf{TB13.}\ Use a left rectangular sum on the five table intervals to estimate $\int_0^{28}f(x)\,dx$: $\begin{array}{c|rrrrrr}x&0&4&10&18&22&28\\\hline f(x)&8&12&6&2&0&3\end{array}$.
\begin{list}{}{\setlength{\leftmargin}{3em}\setlength{\labelwidth}{2.5em}\setlength{\itemsep}{0pt}\setlength{\topsep}{2pt}}
\item[(A)] $130$
\item[(B)] $150$
\item[(C)] $160$
\item[(D)] $175$
\end{list}

\bigskip\noindent\textbf{TB14.}\ Find the value of $x$ at which the function $y=x^2$ reaches its average value on the interval $[0,12]$.
\begin{list}{}{\setlength{\leftmargin}{3em}\setlength{\labelwidth}{2.5em}\setlength{\itemsep}{0pt}\setlength{\topsep}{2pt}}
\item[(A)] $6.351$
\item[(B)] $6.928$
\item[(C)] $7.211$
\item[(D)] $20.78$
\end{list}

\bigskip\noindent\textbf{TB15.}\ The average value of $f(x) = \begin{cases} x^3 & x<3 \\ 5x & x \geq 3 \end{cases}$ on the interval $0 \leq x \leq 6$ is
\begin{list}{}{\setlength{\leftmargin}{3em}\setlength{\labelwidth}{2.5em}\setlength{\itemsep}{0pt}\setlength{\topsep}{2pt}}
\item[(A)] $12.25$
\item[(B)] $13.5$
\item[(C)] $14.625$
\item[(D)] $18$
\end{list}

\end{document}
